\subsection{Las metas de gestión en el régimen tarifario}

El Decreto Legislativo N.\textdegree~1280, que aprueba la Ley del Servicio Universal de Agua Potable y Saneamiento (en adelante, Ley del Servicio Universal), atribuye a la Sunass la competencia para establecer la normativa y los procedimientos aplicables a la regulación económica de los servicios de agua potable y saneamiento.\footnote{Artículo 70 de la Ley del Servicio Universal.} En desarrollo de la Ley, el Reglamento del Decreto Legislativo N.\textdegree~1280, aprobado mediante Decreto Supremo N.\textdegree~009-2024-VIVIENDA (en adelante, Reglamento de la Ley del Servicio Universal), regula el régimen aplicable a las metas de gestión. Conforme con su artículo 168, estas permiten realizar el seguimiento y la evaluación individual del cumplimiento de los resultados del programa de inversiones y de las actividades destinadas al desarrollo o la mejora de la estructura organizacional, entre otras acciones que coadyuven a mejorar la gestión y la prestación de los servicios.\footnote{Párrafo 168.1 del artículo 168 del Reglamento de la Ley del Servicio Universal.}

Sobre esta base normativa, el \textit{Reglamento General de Tarifas de los Servicios de Saneamiento brindados por Empresas Prestadoras}, aprobado mediante Resolución de Consejo Directivo N.\textdegree~028-2021-SUNASS-CD y sus modificatorias (en adelante, Reglamento General de Tarifas), define las metas de gestión como parámetros destinados al seguimiento y la evaluación sistémica del cumplimiento de las inversiones y medidas de mejora programadas para el periodo regulatorio, orientadas a mejorar la gestión y la prestación de los servicios.\footnote{Literal n) del artículo IV del Reglamento General de Tarifas.} Desde la aprobación del Reglamento en 2021, el tratamiento regulatorio de estas metas ha experimentado cambios en tres aspectos relevantes para el análisis: su relación con los incrementos tarifarios, el criterio normativo utilizado para identificar y ordenar los indicadores empleados en su medición y el tratamiento de los indicadores asociados con la ejecución del programa de inversiones. El primero modificó la relación entre el cumplimiento de las metas y la aplicación de los incrementos tarifarios; el segundo consolidó al Sistema de Indicadores e Índices de la Gestión de los Prestadores de los Servicios de Saneamiento (SIIGEPSS) como instrumento para identificar y ordenar los indicadores aplicables; y el tercero delimitó progresivamente el alcance de la verificación financiera del programa de inversiones.

Respecto de los incrementos tarifarios, el diseño aprobado en 2021 permitía que el cumplimiento de las metas de gestión condicionara la aplicación total o parcial de los incrementos tarifarios base programados.\footnote{Véase el apartado 5.8 del Informe N.\textdegree~033-2021-SUNASS-DPN, que sustentó la aprobación del Reglamento General de Tarifas de 2021.} Este vínculo fue eliminado mediante el Decreto Legislativo N.\textdegree~1620, que estableció que dichos incrementos no se encuentran sujetos al cumplimiento de condiciones, sin perjuicio de las sanciones que correspondan por el incumplimiento de las metas de gestión.\footnote{Párrafo 71.3 del artículo 71 de la Ley del Servicio Universal. Esta disposición fue desarrollada en los párrafos 168.3 y 168.4 del artículo 168 del Reglamento de la Ley del Servicio Universal.} Mediante Resolución de Consejo Directivo N.\textdegree~037-2024-SUNASS-CD se adecuó el Reglamento General de Tarifas a este nuevo marco, eliminando la disposición según la cual la aplicación de los incrementos tarifarios base estaba sujeta al cumplimiento de las metas de gestión.\footnote{Párrafo 19.1 del artículo 19 del Reglamento General de Tarifas.}

Una vez desvinculados los incrementos tarifarios base del cumplimiento de condiciones, se mantuvo la clasificación entre metas de gestión base y condicionadas, sustentada en el nivel de certeza asociado a la ejecución de las inversiones. La adecuación aprobada mediante Resolución de Consejo Directivo N.\textdegree~124-2025-SUNASS-CD precisó los criterios para determinar dicho nivel de certeza, considerando especialmente el grado de madurez de las intervenciones y la fase alcanzada dentro del ciclo de inversión.\footnote{El alto nivel de certeza puede sustentarse, entre otros supuestos, en la existencia de un expediente técnico o documento equivalente aprobado o en el grado de madurez alcanzado por el proyecto conforme con la normativa de inversiones.} De esta manera, una intervención puede recibir un tratamiento diferente según su grado de madurez, incluida la posibilidad de asociar etapas posteriores de ejecución u operación a incrementos condicionados.

\begin{table*}[t]
\caption{Indicadores de gestión aplicables a las empresas prestadoras}
\label{tab:grupos_indicadores_ep}
\centering
\scriptsize
\renewcommand{\arraystretch}{1.25}
\setlength{\tabcolsep}{5pt}
\begin{tabularx}{\textwidth}{
@{}
>{\raggedright\arraybackslash}p{0.19\textwidth}
>{\centering\arraybackslash}p{0.11\textwidth}
>{\raggedright\arraybackslash}X
@{}
}
\hline
\textbf{Aspecto} & \textbf{Indicadores} & \textbf{Contenido y ejemplos} \\
\hline
Cobertura & EP-01 a EP-06 & Comprende indicadores de cobertura de agua potable y alcantarillado, así como indicadores referidos a la instalación de nuevas conexiones domiciliarias y piletas. \\
Calidad del servicio & EP-07 a EP-27 & Comprende indicadores de continuidad, presión, turbiedad, cloro residual, interrupciones y reclamos, así como indicadores relacionados con atoros, roturas, renovación de redes y tratamiento de aguas residuales. \\
Sostenibilidad de los servicios & EP-28 a EP-30 & Comprende la relación de trabajo, la ejecución de reservas y la incorporación de procesos de gestión del riesgo de desastres y adaptación al cambio climático. \\
Solvencia económica y financiera & EP-31 a EP-39 & Comprende indicadores de liquidez, endeudamiento, cobertura del servicio de deuda, márgenes operativo, neto y EBITDA, y rendimiento sobre los activos y el patrimonio. \\
Ganancia de eficiencia empresarial & EP-40 a EP-56 & Comprende indicadores como morosidad, costos operativos, agua no facturada, macromedición y micromedición; catastros comercial y técnico; reemplazo de medidores; y recuperación de conexiones inactivas. \\
Cumplimiento del programa de inversiones & EP-57 y EP-57-A & Comprende el avance financiero del programa de inversiones y el avance financiero focalizado de las inversiones y medidas de mejora cuyo impacto no pueda medirse mediante otro indicador. \\
Otros indicadores & EP-58 a EP-60 & Comprende indicadores de buen gobierno corporativo, cumplimiento de obligaciones y sobrecobertura del subsidio. \\
\hline
\end{tabularx}
\begin{minipage}{0.97\textwidth}
\footnotesize
\textit{Nota:} La tabla presenta la clasificación por aspectos establecida en el SIIGEPSS y comprende únicamente los indicadores de gestión aplicables a las empresas prestadoras. No incluye los índices de gestión definidos por dicho sistema. El indicador EP-57-A fue incorporado mediante Resolución de Consejo Directivo N.\textdegree~124-2025-SUNASS-CD; su alcance y relación con el EP-57 se desarrollan en el texto.
\end{minipage}
\end{table*}

En línea con esta clasificación, el Reglamento de la Ley del Servicio Universal incorporó el Programa de Ejecución de Metas Físico-Financieras (PEMFF) como instrumento de programación para el cumplimiento de las metas de gestión. El PEMFF forma parte del estudio tarifario e identifica los supuestos utilizados para formular las metas, los plazos estimados del ciclo de inversión, las alternativas de implementación y los principales riesgos y medidas de contingencia.\footnote{Artículo 169 del Reglamento de la Ley del Servicio Universal.} La adecuación del Reglamento General de Tarifas de 2025 desarrolló este instrumento para las metas de gestión base y estableció un seguimiento periódico de la ejecución de las inversiones.\footnote{Artículos 15-A y 15-B del Reglamento General de Tarifas.} El PEMFF permite relacionar las metas aprobadas con los supuestos que sustentaron su programación y con las circunstancias que pueden afectar su ejecución.

En cuanto al ordenamiento normativo de los indicadores, el cambio se concentró en el criterio utilizado para identificarlos y clasificarlos. En la versión del Reglamento General de Tarifas aprobada en 2021, el artículo 13 organizaba las metas de gestión según seis objetivos regulatorios: cumplimiento del programa de inversiones, calidad del servicio, solvencia económica y financiera de la empresa prestadora, sostenibilidad de los servicios, ganancia de eficiencia empresarial y cobertura o cierre de brechas. Esta clasificación ordenaba las metas según el objetivo regulatorio al que contribuían y permitía vincularlas con la finalidad regulatoria correspondiente. El Reglamento disponía, además, que los indicadores utilizados para medirlas se encontraban establecidos en el sistema de indicadores aprobado por la Sunass o en aquel que lo reemplazara.\footnote{Redacción del artículo 13 del Reglamento General de Tarifas anterior a su modificación mediante Resolución de Consejo Directivo N.\textdegree~124-2025-SUNASS-CD.} Ese mismo año, mediante Resolución de Consejo Directivo N.\textdegree~063-2021-SUNASS-CD, se aprobó el SIIGEPSS, que organizó los indicadores aplicables a las EP en aspectos que guardaban correspondencia con los objetivos regulatorios previstos en el Reglamento General de Tarifas.

La adecuación del artículo 13 aprobada mediante Resolución de Consejo Directivo N.\textdegree~124-2025-SUNASS-CD dejó de enumerar los seis objetivos regulatorios y pasó a establecer que las metas de gestión tienen como objetivo la mejora del servicio. Al mismo tiempo, dispuso que los indicadores utilizados para medir su cumplimiento se encuentran comprendidos en el SIIGEPSS.\footnote{Véase el análisis de la modificación del artículo 13 contenido en el Informe N.\textdegree~164-2025-SUNASS-DPN, que sustenta la Resolución de Consejo Directivo N.\textdegree~124-2025-SUNASS-CD.} Con ello, la clasificación por aspectos prevista en dicho sistema pasó a constituir el criterio normativo para identificar y ordenar los indicadores susceptibles de utilizarse en la determinación de las metas de gestión. La Tabla~\ref{tab:grupos_indicadores_ep} presenta esta clasificación.

De otro lado, en relación con la medición del cumplimiento del programa de inversiones, antes de la aprobación del Reglamento General de Tarifas de 2021 las metas establecidas en los estudios tarifarios ya comprendían distintas dimensiones de la gestión y de la prestación, pero no incluían una meta específica destinada a evaluar el programa de inversiones en su conjunto. Entre las 41 EP que culminaron sus periodos regulatorios entre 2014 y 2019 se establecieron metas asociadas con la calidad del servicio, la eficiencia operativa y comercial, la implementación de instrumentos de gestión y la ampliación o mejora de la infraestructura. Su frecuencia de incorporación difería considerablemente, pero ninguna evaluaba específicamente el cumplimiento del programa de inversiones.\footnote{La micromedición fue incluida en 34 empresas, mientras que las conexiones inactivas fueron consideradas únicamente en dos. El indicador de micromedición comprendía la instalación de nuevos medidores y la renovación de medidores. Véase el apartado 4.6.3 del Informe N.\textdegree~033-2021-SUNASS-DPN.}

El diagnóstico regulatorio que sustentó el Reglamento General de Tarifas de 2021 advirtió la existencia de problemas de ejecución en las EP y, al mismo tiempo, la ausencia de una meta específica que evaluara el cumplimiento de sus programas de inversiones. En ese contexto, el Reglamento incorporó el \textit{Cumplimiento del programa de inversiones} como uno de los objetivos regulatorios de las metas de gestión. El SIIGEPSS operativizó este objetivo mediante el indicador EP-57, denominado \textit{Porcentaje de avance financiero del programa de inversiones de la EP}.

En su formulación original, el EP-57 fue diseñado para medir la ejecución financiera acumulada del programa de inversiones y, por regla general, comprendía el conjunto de inversiones reconocidas en el estudio tarifario. Este alcance podía incluir intervenciones vinculadas, a su vez, con otras metas. Era el caso, por ejemplo, de la instalación de \textit{data loggers} asociada con la meta de continuidad del servicio o de la instalación de medidores asociada con la meta de micromedición. Una misma intervención podía formar parte del indicador de avance financiero y estar simultáneamente vinculada con otra meta referida a su ejecución física o a un resultado esperado en la prestación.

La meta de avance financiero alcanzó una presencia amplia en los estudios tarifarios aprobados durante los años siguientes. Entre 2021 y 2024, 41 de los 48 estudios analizados incorporaron esta meta, lo que equivale al 85,4\,\% y representa una de las mayores frecuencias observadas durante el periodo. La Tabla~\ref{tab:metas_gestion_estudios_recientes} compara esta frecuencia con la correspondiente a las demás metas aprobadas.

\begin{table}[!t]
\centering
\caption{Metas de gestión aprobadas en los estudios tarifarios, 2021--2024}
\label{tab:metas_gestion_estudios_recientes}
\scriptsize
\setlength{\tabcolsep}{2.5pt}
\renewcommand{\arraystretch}{1.08}

\begin{tabularx}{\columnwidth}{
>{\centering\arraybackslash}p{0.45cm}
>{\raggedright\arraybackslash}X
>{\centering\arraybackslash}p{1.35cm}
}
\hline

\parbox[c][0.8cm][c]{0.45cm}{\centering\textbf{N.°}}
&
\parbox[c][0.8cm][c]{\linewidth}{\centering\textbf{Indicador de la meta de gestión}}
&
\parbox[c][0.8cm][c]{1.35cm}{\centering\textbf{N.° de\\empresas\\prestadoras}}
\\

\hline
1  & Continuidad & 48 \\
2  & Relación de trabajo & 48 \\
3  & Catastro comercial & 45 \\
4  & Porcentaje de ejecución de la reserva de MRSE & 45 \\
5  & Presión & 42 \\
6  & Porcentaje de avance financiero del programa de inversiones de la EP\textsuperscript{b} & 41 \\
7  & Catastro técnico & 40 \\
8  & Reemplazo de medidores & 37 \\
9  & Agua no facturada & 27 \\
10 & Porcentaje de ejecución de la reserva para GRD y ACC & 26 \\
11 & Micromedición & 25 \\
12 & Porcentaje de ejecución de la reserva para PCC y PAS & 18 \\
13 & Porcentaje de ejecución de la reserva para la GRD & 17 \\
14 & Porcentaje de ejecución de la reserva para la implementación del PCC & 15 \\
15 & Instalación de medidores & 12 \\
16 & Recuperación de conexiones inactivas del servicio de agua potable & 10 \\
17 & Otros indicadores\textsuperscript{a} & 25 \\
\hline
\end{tabularx}

\vspace{0.15cm}
\begin{minipage}{\columnwidth}
\scriptsize
\textit{Nota:} La tabla comprende 48 estudios tarifarios aprobados entre 2021 y 2024 y presenta individualmente los indicadores asignados a diez o más empresas prestadoras. \textsuperscript{a}Agrupa once indicadores asignados a menos de diez empresas prestadoras. \textsuperscript{b}En dos EP, la meta financiera tuvo desde su aprobación un alcance delimitado; estos casos y la posterior revisión excepcional de EPS GRAU S.A. se desarrollan en el texto.\\
\textit{Fuente:} Estudios tarifarios de la Sunass.\\
\end{minipage}
\end{table}

En ese marco, el Reglamento de la Ley del Servicio Universal, aprobado en 2024, introdujo una restricción expresa al uso de la medición financiera. Su artículo 168 establece que el porcentaje de avance financiero del programa de inversiones, o el indicador que haga sus veces, se aprueba excepcionalmente cuando no pueda aplicarse otro indicador de gestión o de calidad del servicio. Esta regla guarda relación con la exigencia, prevista en el mismo artículo, de evitar duplicidades en la evaluación de las metas de gestión.\footnote{Párrafos 168.1 y 168.2 del artículo 168 del Reglamento de la Ley del Servicio Universal.} Así, la medición financiera queda reservada a aquellos supuestos en los que el cumplimiento no pueda evaluarse mediante otro indicador aplicable.

En aplicación de estas reglas, la Resolución de Consejo Directivo N.\textdegree~124-2025-SUNASS-CD incorporó al SIIGEPSS el indicador EP-57-A, denominado \textit{Porcentaje de avance financiero focalizado del programa de inversiones y medidas de mejora de la EP}, y precisó que su utilización procede únicamente cuando el impacto de las inversiones o medidas de mejora no puede medirse mediante un indicador distinto.\footnote{Párrafo 13.2 del artículo 13 del Reglamento General de Tarifas y párrafo 168.2 del artículo 168 del Reglamento de la Ley del Servicio Universal.} El indicador permite, por ejemplo, evaluar financieramente etapas como la elaboración de expedientes técnicos o medidas de mejora cuyo efecto todavía no se manifiesta directamente en un indicador de gestión o de calidad del servicio.

Cabe precisar que la incorporación del EP-57-A no eliminó el EP-57 del SIIGEPSS. Tampoco modificó automáticamente las metas aprobadas con anterioridad, que continúan siendo exigibles conforme con los estudios y resoluciones tarifarias que las establecieron, mientras no sean modificadas mediante el procedimiento correspondiente. En la práctica, pueden coexistir metas de avance financiero con alcances diferentes: algunas referidas al programa de inversiones en su conjunto y otras circunscritas a inversiones y medidas de mejora que no son evaluadas mediante indicadores distintos.

En síntesis, la evolución descrita deja tres antecedentes especialmente relevantes para el análisis que sigue. Primero, el cumplimiento de las metas dejó de condicionar la aplicación de los incrementos tarifarios base. Segundo, el ordenamiento normativo pasó de una clasificación explícita de las metas por objetivos regulatorios a una identificación y ordenación de los indicadores a partir del SIIGEPSS. Tercero, la medición financiera del programa de inversiones fue delimitándose frente a aquellos casos en los que el cumplimiento puede evaluarse mediante otros indicadores.